% Libro 4 · Mundo 6 · Lectura de dos textos (app: pu-dos-textos)
\mundo{Lectura de dos textos}{pu-dos-textos}{A veces el examen presenta dos textos sobre el mismo tema y pregunta por la
relación entre ellos: en qué coinciden, en qué se oponen, qué diría un autor de la idea del otro. Es la tarea más completa
de la prueba, porque exige leer cada texto por separado y, además, ponerlos a conversar.}{%
\item \textbf{6.1} Coincidencias y diferencias
\item \textbf{6.2} Qué diría el otro autor}

\leccion{pu-dt-1}{Coincidencias y diferencias}{Juego de la tabla de dos columnas.}

\begin{reto}
El texto 1 dice que los videojuegos aíslan a los jóvenes; el texto 2, que pueden crear comunidades. ¿En qué coinciden,
probablemente, los dos autores?
\opcionesl{En que los videojuegos son siempre dañinos.}{En que los videojuegos influyen en la vida social de los jóvenes.}{En que hay que prohibirlos.}{En nada.}
\end{reto}

\tecnica{La técnica: la tabla de dos columnas}
Con dos textos, la herramienta es una tabla: una columna por texto y una fila por aspecto (tema, tesis, argumentos, tono,
propuesta). Al llenarla se ve dónde coinciden y dónde se separan.

\begin{center}\small
\begin{tabular}{@{}l|p{40mm}|p{40mm}@{}}
\toprule
 & \textbf{Texto 1} & \textbf{Texto 2}\\\midrule
Tema & & \\
Tesis & & \\
Argumentos & & \\
Tono & & \\\bottomrule
\end{tabular}
\end{center}

\subsubsection{Tipos de relación}
Dos textos pueden \emph{coincidir} (dicen lo mismo con distintas razones), \emph{oponerse} (defienden tesis contrarias),
\emph{complementarse} (uno aporta lo que al otro le falta) o tratar \emph{aspectos distintos} del mismo tema. Ojo: dos
textos que se oponen en la tesis pueden coincidir en el diagnóstico.

\begin{regla}
\begin{itemize}
  \item Lee cada texto por separado y escribe su tesis en una línea.
  \item Llena la tabla: tema, tesis, argumentos, tono.
  \item En las preguntas de coincidencia, busca el punto común más preciso, no el más amplio.
\end{itemize}
\end{regla}

\begin{texto}{Texto 1 · La tarea tiene los días contados}{Columna · texto original escrito para este libro}
\pa Las tareas para la casa deberían reducirse al mínimo. Los estudiantes pasan siete u ocho horas en el colegio; agregar
tres horas más de deberes les quita tiempo para descansar, hacer deporte y estar con su familia.

\pa Además, la tarea aumenta la desigualdad: el estudiante que tiene en casa a alguien que le ayuda avanza, y el que no,
se atrasa. El colegio termina premiando las condiciones del hogar, no el esfuerzo del estudiante.
\end{texto}

\begin{texto}{Texto 2 · Defensa de la tarea bien hecha}{Columna · texto original escrito para este libro}
\pa Eliminar la tarea sería un error. Hay aprendizajes que solo se consolidan con práctica individual, a solas con el
problema, sin el profesor al lado. Esa práctica no cabe en las horas de clase.

\pa Reconozco, sin embargo, que mucha tarea es inútil: copias, cuestionarios repetitivos, trabajos que terminan
haciendo los padres. El problema no es la tarea, sino la mala tarea. Pocas, breves y que el estudiante pueda hacer solo:
esas son las tareas que valen.
\end{texto}

\subsection{Preguntas comentadas}
\begin{comentada}{¿En qué coinciden los dos autores?}
\mal{A}{En que la tarea debe eliminarse por completo.}{El texto 2 lo rechaza; el texto 1 dice «reducirse al mínimo», no
eliminarse.}
\bien{B}{En que hay tareas que no deberían depender de la ayuda de los padres.}{Texto 1: la ayuda en casa crea
desigualdad; texto 2: critica los trabajos que «terminan haciendo los padres» y pide tareas que el estudiante haga solo.}
\mal{C}{En que la práctica individual es indispensable.}{Solo lo dice el texto 2.}
\mal{D}{En que los estudiantes deben pasar más horas en el colegio.}{\emph{No lo dice} ninguno.}
\end{comentada}

\taller
\piso{Calienta}{la tabla}
\begin{ejercicios}
\item Llena la tabla de dos columnas con tema, tesis y argumentos. \sol{Tema: la tarea escolar. Tesis 1: reducirla al mínimo; tesis 2: mantenerla, pero buena. Argumentos 1: tiempo de descanso, desigualdad; argumentos 2: la práctica individual consolida, la mala tarea es el problema.}
\end{ejercicios}
\piso{Como en el examen}{opción múltiple}
\begin{ejercicios}
\item La relación entre las tesis de los dos textos es
\opcionesl{de identidad: dicen lo mismo.}{de oposición parcial: uno quiere reducirla al mínimo y el otro mantenerla con condiciones.}{de ejemplo: el 2 ejemplifica el 1.}{no tienen relación.}
\sol{B.}
\item ¿Qué argumento del texto 1 tiene una respuesta parcial en el texto 2?
\opcionesl{El tiempo para hacer deporte.}{La desigualdad causada por la ayuda en casa.}{Las horas en el colegio.}{La familia.}
\sol{B. El texto 2 pide tareas que el estudiante pueda hacer solo.}
\item El tono de ambos textos es
\opcionesd{argumentativo}{narrativo}{humorístico}{instructivo}
\sol{A.}
\end{ejercicios}

\begin{retoresuelto}
Aunque llegan a conclusiones opuestas, ambos suponen que los videojuegos afectan la vida social de los jóvenes (para mal o
para bien). Ese es el punto común preciso. \textbf{Opción B.}
\end{retoresuelto}
\bitacora

\leccion{pu-dt-2}{Qué diría el otro autor}{Juego de responder con la voz del otro texto.}

\begin{reto}
Si el autor del texto 2 (defensa de la tarea bien hecha) leyera la frase «hay que eliminar toda tarea», lo más probable es
que
\opcionesl{estuviera de acuerdo.}{la rechazara, porque cree que la práctica individual es necesaria.}{no opinara.}{propusiera más horas de clase.}
\end{reto}

\tecnica{La técnica: ponerse en los zapatos del autor}
Para responder «¿qué diría el autor A sobre la idea B?», no se usa la opinión propia: se usa la \emph{lógica del autor}.
Se toma su tesis y sus argumentos y se aplican a la idea nueva. Si la idea contradice su tesis, la rechazaría; si la
apoya, la aceptaría; si coincide en parte, la matizaría.

\begin{regla}
\begin{enumerate}
  \item Escribe la tesis y los valores del autor (qué le importa).
  \item Confronta la idea nueva con esos valores.
  \item Elige la respuesta coherente con todo el texto, no con una frase suelta.
\end{enumerate}
\end{regla}

\begin{texto}{Texto 1 · El celular en el aula}{Columna · texto original escrito para este libro}
\pa Prohibir el celular en los colegios es la única manera de recuperar la atención. Un estudiante con el teléfono en el
bolsillo está, en realidad, en dos lugares a la vez: en la clase y en las notificaciones. Ningún profesor puede competir
con eso. Los colegios que lo prohibieron reportan recreos más conversados y clases más tranquilas.
\end{texto}

\begin{texto}{Texto 2 · Enseñar a usarlo}{Columna · texto original escrito para este libro}
\pa Prohibir el celular es renunciar a enseñar a usarlo. Los jóvenes van a vivir con estos aparatos toda la vida; el
colegio debería ser el lugar donde aprendan a apagarlos cuando hace falta, a buscar información confiable y a no dejarse
manejar por las notificaciones. Una norma clara de uso enseña más que una prohibición.
\end{texto}

\subsection{Preguntas comentadas}
\begin{comentada}{¿Qué diría el autor del texto 2 sobre los recreos «más conversados» que menciona el texto 1?}
\mal{A}{Que demuestran que hay que prohibir el celular.}{Es la conclusión del texto 1, no la del 2.}
\bien{B}{Que el mismo resultado podría lograrse enseñando cuándo apagar el celular, sin prohibirlo.}{Coherente con su
tesis: aprender a apagarlos «cuando hace falta».}
\mal{C}{Que conversar en el recreo es inútil.}{\emph{No lo dice} ni se sigue.}
\mal{D}{Que los profesores deben competir con las notificaciones.}{\emph{Dice otra cosa}.}
\end{comentada}

\taller
\piso{Como en el examen}{opción múltiple}
\begin{ejercicios}
\item El autor del texto 1 probablemente respondería a la propuesta de «una norma clara de uso» que
\opcionesl{es la mejor solución.}{no basta, porque el celular en el bolsillo ya divide la atención.}{los celulares deben ser gratis.}{hay que enseñar a programar.}
\sol{B.}
\item ¿En qué coinciden ambos autores?
\opcionesl{En prohibir el celular.}{En que el celular puede afectar la atención y el colegio debe hacer algo al respecto.}{En que el celular es inofensivo.}{En que los profesores no saben usarlo.}
\sol{B.}
\item ¿Qué valor es más importante para el autor del texto 2?
\opcionesd{la educación para la autonomía}{la disciplina estricta}{el silencio}{la tecnología por sí misma}
\sol{A.}
\end{ejercicios}

\begin{retoresuelto}
El texto 2 defiende la práctica individual como necesaria; eliminar toda tarea la contradice. Lo rechazaría, aunque
aceptaría reducir la mala tarea. \textbf{Opción B.}
\end{retoresuelto}
\bitacora

\evaluacion
\begin{texto}{Texto 1 · Turismo que salva}{Columna · texto original escrito para este libro}
\pa El turismo puede ser la mejor herramienta de conservación que tenemos. Un pueblo que vive de mostrar sus aves, sus
cascadas o sus bosques tiene una razón económica para cuidarlos. Donde antes se talaba para ganadería, hoy hay guías
locales que ganan más mostrando un colibrí que vendiendo madera.
\end{texto}
\begin{texto}{Texto 2 · Turismo que agota}{Columna · texto original escrito para este libro}
\pa El turismo masivo está matando los lugares que dice amar. Senderos erosionados, basura en las cascadas, arriendos
impagables para los habitantes. Cuando un pueblo se vuelve «de moda», los primeros en irse son los que lo hicieron
especial. Sin límites de visitantes, el turismo consume lo que vende.
\end{texto}
\begin{ejercicios}
\item La tesis del texto 1 es que el turismo
\opcionesl{destruye la naturaleza.}{puede motivar la conservación al darle valor económico.}{debe prohibirse.}{solo beneficia a los extranjeros.}
\sol{B. (6.1)}
\item Los dos textos coinciden en que
\opcionesl{el turismo tiene efectos importantes sobre los lugares que visita.}{el turismo siempre es bueno.}{hay que talar menos.}{los arriendos deben subir.}
\sol{A. (6.1)}
\item El autor del texto 2 respondería al ejemplo del colibrí del texto 1 que
\opcionesl{prueba que el turismo no tiene riesgos.}{ese beneficio puede perderse si el turismo crece sin límites.}{los colibríes no existen.}{los guías ganan poco.}
\sol{B. (6.2)}
\item ¿Qué propuesta podría aceptar el autor de cada texto?
\opcionesl{Prohibir todo turismo.}{Un turismo con límites de visitantes que beneficie a los habitantes.}{Talar para ganadería.}{Subir los arriendos.}
\sol{B. Al texto 1 le sirve para conservar; al 2, porque pone límites. (6.2)}
\end{ejercicios}
